\begin{abstract}
O desenvolvimento de aplicações que utilizam Inteligência Artificial Generativa (Gen AI) tem crescido de forma exponencial. No entanto, a ausência de um protocolo comum na comunicação com Modelos de Linguagem de Grande Escala (LLMs) obriga estas aplicações a integrarem diferentes interfaces de acesso e implementações particulares para a invocação de ferramentas externas. Atualmente cada fornecedor (como \textbf{OpenAI} \cite{openai_api}, \textbf{Anthropic} \cite{anthropic_api} e \textbf{Gemini} \cite{google_gemini_api}) impõe as suas próprias estruturas de dados, ciclos de vida de \textit{streaming} e padrões para a invocação de ferramentas (\textit{tool calling})\cite{ibm_tool_calling}\cite{toolCallingCycle}. Para resolver esta fragmentação tecnológica, este projeto propõe o desenvolvimento do \textbf{OmniAI SDK}, um \textit{Software Development Kit} construído em \textbf{Kotlin Multiplatform} (\textbf{KMP}) \cite{kmp_ref}. O objetivo principal desta biblioteca é fornecer a infraestrutura base necessária para instanciar e configurar de forma rápida uma \textbf{OmniAI Gateway} (um ponto de entrada centralizado dedicado ao acesso de inferência, segurança, observabilidade e outros). 
O \textbf{OmniAI SDK} é baseado nos princípios da Arquitetura Hexagonal (\textit{Ports and Adapters}) \cite{hexagonal_ports}, este isola a lógica de domínio e suporta a comunicação bidirecional compatível com os padrões das APIs de referência, nomeadamente \textbf{OpenAI} \cite{openai_api}, \textbf{Anthropic} \cite{anthropic_api} e \textbf{Gemini} \cite{google_gemini_api}. O sistema implementa uma \textit{pipeline} extensível e configurável, 
inspirada no modelo de filtros e cadeia de filtros dos \textbf{Servlets} \cite{java_servlet_filter}\cite{java_servlet_filterchain}, que executa \textit{interceptors} pré-construídos para autorização, 
autenticação, métricas, \textit{rate limiting}, \textit{circuit breaker} e \textit{fallback}, podendo ser adicionados mais \textit{interceptors} customizados pelo utilizador do \textbf{OmniAI SDK}. 
Adicionalmente, o \textit{Software Development Kit} prepara um módulo \textbf{MCP Broker} para evoluir a \textbf{OmniAI Gateway} de uma \textit{gateway} de inferência e torná-la numa \textit{AI Gateway} completa, adicionando a gestão de ferramentas externas à \textbf{OmniAI Gateway} utilizando o \textbf{Model Context Protocol (MCP)} \cite{mcp_protocol}. Distribuído de forma multiplataforma para ambientes \textbf{JVM}, 
via \textbf{Maven}, e \textbf{Node.js}/\textbf{TypeScript}, via \textbf{NPM}, o \textbf{OmniAI SDK} foi validado através da construção de uma \textbf{OmniAI Gateway}, 
um \textit{Proof of Concept} que utiliza a biblioteca para demonstrar a criação de uma \textit{gateway} real. A avaliação demonstra que a solução 
minimiza o acoplamento tecnológico e o esforço de integração, garantindo resiliência, escalabilidade e centralização na utilização de clientes de IA, como o \textbf{Claude Code} \cite{claude_code}, entre outros.
\end{abstract}

\renewcommand{\abstractname}{Abstract}
\begin{abstract}
The development of applications leveraging Generative Artificial Intelligence (Gen AI) has grown exponentially. However, the lack of a common protocol for communicating with Large Language Models (LLMs) forces these applications to integrate different access interfaces and custom implementations for invoking external tools. Currently, each provider (such as \textbf{OpenAI} \cite{openai_api}, \textbf{Anthropic} \cite{anthropic_api}, and \textbf{Gemini} \cite{google_gemini_api}) imposes its own data structures, \textit{streaming} lifecycles, and patterns for \textit{tool calling} \cite{ibm_tool_calling}\cite{toolCallingCycle}. To address this technological fragmentation, this project proposes the development of the \textbf{OmniAI SDK}, a \textit{Software Development Kit} built in \textbf{Kotlin Multiplatform} (\textbf{KMP}) \cite{kmp_ref}. The main objective of this library is to provide the necessary foundational infrastructure to quickly instantiate and configure an \textbf{OmniAI Gateway} (a centralized entry point dedicated to inference access, security, observability, and others).

Based on the principles of Hexagonal Architecture (\textit{Ports and Adapters}) \cite{hexagonal_ports}, the \textbf{OmniAI SDK} isolates the domain logic and supports bidirectional communication compatible with the standards of reference APIs, namely \textbf{OpenAI} \cite{openai_api}, \textbf{Anthropic} \cite{anthropic_api}, and \textbf{Gemini} \cite{google_gemini_api}. The system implements an extensible and configurable \textit{pipeline}, inspired by the \textbf{Servlets} filter and filter chain model \cite{java_servlet_filter}\cite{java_servlet_filterchain}, which executes pre-built \textit{interceptors} for authorization, authentication, metrics, \textit{rate limiting}, \textit{circuit breaking}, and \textit{fallback}. Additionally, custom \textit{interceptors} can be added by the \textbf{OmniAI SDK} user.

Furthermore, the \textit{Software Development Kit} includes an \textbf{MCP Broker} module to evolve the \textbf{OmniAI Gateway} from an inference \textit{gateway} into a full-fledged \textit{AI Gateway}, by adding external tool management capabilities using the \textbf{Model Context Protocol (MCP)} \cite{mcp_protocol}. Distributed as a multiplatform solution for \textbf{JVM} environments via \textbf{Maven}, and for \textbf{Node.js}/\textbf{TypeScript} via \textbf{NPM}, the \textbf{OmniAI SDK} was validated through the construction of an \textbf{OmniAI Gateway}, a \textit{Proof of Concept} that utilizes the library to demonstrate the creation of a real-world \textit{gateway}. The evaluation demonstrates that the solution minimizes technological coupling and integration effort, ensuring resilience, scalability, and centralization when using AI clients, such as \textbf{Claude Code} \cite{claude_code}, among others.
\end{abstract}

\tableofcontents
\listoffigures
\newpage